\section{State of the Art and Baselines}
\label{sec:app1_baselines_sota}

\paragraph{State of the Art}
\citet{ethayarajh-2019-contextual} showed that contextual representations of BERT, ELMo, and GPT-2 are anisotropic: random words have high average cosine similarity, and this grows with depth. That work introduced the average random cosine as an anisotropy baseline and maximum explainable variance (MEV) as the share of variance on the first principal component. \citet{rudman-etal-2022-isoscore} showed that average cosine and partition-based scores do not measure uniform use of dimensions and proposed IsoScore, which is 1 when variance is spread equally over all $d$ dimensions and 0 when it lies on one direction. They also used the maximum likelihood intrinsic dimension (ID) of \citet{levina-bickel-2004-maximum}. Both works study English text only.

\paragraph{Competitive Baselines}
Our baselines are the metrics and the models. We report average random cosine and MEV \citep{ethayarajh-2019-contextual} next to IsoScore and ID \citep{rudman-etal-2022-isoscore}, so results stay comparable to both papers. Within a model family, the base model is the baseline for the embedding model: Gemma-3-1B-PT for EmbeddingGemma-300M, and Qwen3.5-0.8B-Base for Qwen3-Embedding-0.6B. English (\texttt{eng\_Latn}) is the language baseline for the 22 Indic languages.

\paragraph{Baseline Availability}
The official IsoScore implementation is available as the \texttt{IsoScore} Python package, and \citet{rudman-etal-2022-isoscore} compute ID with \texttt{scikit-dimension}. Both are written in NumPy and run on the CPU. With 16{,}000 points per layer, up to 30 layers, 3 views, 4 models, and 23 languages, this is slow and requires copying every hidden state off the GPU. We therefore ported IsoScore (steps 2 to 7 of the paper), MEV, average random cosine, and the ID estimator to PyTorch so that they run on the GPU next to inference. A unit test checks that our IsoScore matches the official package within $10^{-6}$ on a rotated, shifted, anisotropic Gaussian. Further tests check the expected values on isotropic and rank-one data and recover the dimension of a known linear subspace.
